% Table for ch8-45 -- Partition Tolerance (S10.3), a classification
% (protocol x behavior x what resumes it).
%
% Reader question: when the network partitions, which protocol keeps
%   working, which stalls, and which waits?
% Claim: partition tolerance is per-protocol, not a global property; each
%   of the three governing protocols degrades differently and resumes by a
%   different mechanism.
% Evidence: S:fh-partition's own three one-sentence claims.
% Grammar rejected: a figure would draw one partition event and three
%   reactions as if simultaneous; the claim is that each protocol answers
%   independently, which a table's per-row independence states directly.
% Five-second test: which protocol keeps making progress inside a
%   partition, and which one just waits?
\begin{table}[H]
\centering
\caption{Federation runs no consensus, so partition tolerance is decided
protocol by protocol. Transfer stalls outright; revocation keeps working
inside each side of the split; settlement holds its position until the
split heals. [internal]}
\label{tab:fh-partition}
\small
\renewcommand{\arraystretch}{1.2}
\begin{tabularx}{\linewidth}{@{}l >{\raggedright\arraybackslash}X >{\raggedright\arraybackslash}X@{}}
\toprule
\textbf{protocol} & \textbf{behavior under partition} & \textbf{what resumes it} \\
\midrule
transfer & cannot proceed: cross-realm authentication needs both daemons to reach the witness log & the partition heals and both sides reach the log \\
revocation & continues inside each partition; convergence bounded by that partition's own size & the two partitions merge their informed sets on reconnection \\
settlement & parks at the escrow, holding the bond in place & the partition heals \\
\bottomrule
\end{tabularx}
\end{table}
